% Integralen. Seurat: duizenden kleine stipjes vormen samen een geheel, zoals kleine stukjes bij een integraal.
% Motief: de band is de oppervlakte onder een grafiek, met rechts de rechthoeken van een bovensom,
% die in het schilderij steken. De stipjes onder de kromme verwijzen naar Seurat.
\def\omslagbestand{omslag-grande-jatte.jpg}
\def\omslagbron{Georges Seurat, \textit{Een zondagnamiddag op het eiland La Grande Jatte}, 1884-1886 (publiek domein)}
\def\omslagrand{10.75 + 3.7*(\x/21)^2.6 + 0.4*sin(0.7*\x r)}
\def\omslagmotief{%
  % bovensom op [a, b], op de as boven het logo
  \foreach \links in {14,...,19}{%
    \pgfmathsetmacro{\x}{\links+1}%
    \pgfmathsetmacro{\hoogte}{\omslagrand}%
    \fill[white,fill opacity=.17] (\links,3.7) rectangle (\links+1,\hoogte);
    \draw[white,line width=.5pt,draw opacity=.85] (\links,3.7) rectangle (\links+1,\hoogte);
  }
  \draw[white,line width=.8pt,-{Stealth[length=2.4mm]}] (13.3,3.7) -- (20.6,3.7);
  \foreach \plek/\naam in {14/a,20/b}{%
    \draw[white,line width=.8pt] (\plek,3.6) -- (\plek,3.8);
    \node[text=white,font=\small,inner sep=1pt,anchor=north] at (\plek,3.6) {$\naam$};
  }
  % stipjes onder de kromme: kleiner naarmate ze dieper staan
  \foreach \diepte/\straal in {0.32/0.075,0.62/0.055,0.92/0.035}{%
    \foreach \x in {0.25,0.75,...,13.3}{%
      \fill[fel] (\x,{\omslagrand-\diepte}) circle (\straal);
    }
  }
  \randlijn
}
